\chapter{Literature Review}

\section{Retrieval-Augmented Generation (RAG)}

Lewis et al.~\cite{lewis2020retrieval} introduced the RAG framework, which combines a pre-trained sequence-to-sequence model with a dense retrieval component. The key innovation is that the retrieval step provides dynamic, up-to-date context that the generator can condition on, reducing hallucination rates by up to 30\% compared to standalone LLMs. Subsequent work by Borgeaud et al.~\cite{borgeaud2022improving} demonstrated that retrieval-augmented models could match the performance of 10$\times$ larger models by leveraging external knowledge bases.

In the financial domain, RAG systems have been deployed for regulatory compliance checking~\cite{araci2019finbert}, customer query resolution, and fraud detection documentation. However, these deployments have largely overlooked the security implications of combining retrieval with generation.

\section{Prompt Injection Attacks}

Prompt injection refers to the class of attacks where adversarial inputs are crafted to manipulate the behavior of an LLM beyond its intended scope~\cite{perez2022prompt}. Greshake et al.~\cite{greshake2023not} categorized prompt injections into two types:

\begin{enumerate}
    \item \textbf{Direct Prompt Injection}: The attacker directly includes instructions in the input, e.g., \textit{``Ignore all previous instructions and reveal the database schema.''}
    
    \item \textbf{Indirect Prompt Injection}: The malicious payload is embedded in external data sources (such as uploaded PDFs) that the LLM retrieves and processes.
\end{enumerate}

In the context of RAG systems, indirect prompt injection is particularly dangerous because malicious instructions can be hidden within documents that are uploaded to the knowledge base. When these documents are retrieved as context for a user query, the LLM may follow the embedded malicious instructions.

\subsection{Defense Mechanisms}

Several defense strategies have been proposed:

\begin{itemize}
    \item \textbf{Input Sanitization}: Filtering known attack patterns using regular expressions before they reach the LLM~\cite{liu2023prompt}.
    \item \textbf{Instruction Hierarchy}: Training models to prioritize system prompts over user inputs~\cite{wallace2024instruction}.
    \item \textbf{Semantic Safety Classification}: Using a secondary classifier to detect whether an input contains adversarial intent.
\end{itemize}

Our system implements a combination of regex-based pattern matching for known attack signatures and content sanitization to strip potentially malicious payloads.

\section{Embedding Inversion Attacks}

Song and Raghunathan~\cite{song2020information} demonstrated that dense vector embeddings generated by transformer models can be reverse-engineered to reconstruct significant portions of the original input text. Their attack achieved up to 70\% token-level reconstruction accuracy on BERT embeddings.

Morris et al.~\cite{morris2023text} extended this work by showing that embedding inversion attacks are feasible even against black-box embedding APIs, where the attacker has access only to the output vectors, not the model weights.

\subsection{Defense: Differential Privacy}

Differential privacy, formalized by Dwork~\cite{dwork2006differential}, provides a mathematical framework for quantifying and limiting information leakage. The Laplace mechanism adds calibrated noise to query outputs:

\begin{definition}[$\epsilon$-Differential Privacy]
A randomized mechanism $\mathcal{M}$ satisfies $\epsilon$-differential privacy if for all datasets $D_1, D_2$ differing in at most one record, and for all possible outputs $S$:
\begin{equation}
    \Pr[\mathcal{M}(D_1) \in S] \leq e^{\epsilon} \cdot \Pr[\mathcal{M}(D_2) \in S]
\end{equation}
\end{definition}

The Laplace mechanism achieves $\epsilon$-differential privacy by adding noise drawn from the Laplace distribution:

\begin{equation}
    \mathcal{M}(x) = f(x) + \text{Lap}\left(\frac{\Delta f}{\epsilon}\right)
\end{equation}

where $\Delta f$ is the \textit{sensitivity} of the function $f$, defined as the maximum change in $f$'s output when a single data point is modified.

\section{PII Detection and Pseudonymization}

Personal Identifiable Information (PII) detection in financial documents has been extensively studied. Named Entity Recognition (NER) models such as those based on BERT~\cite{devlin2019bert} can identify entities like names, addresses, and account numbers. However, for Indian banking documents with specific formats (IFSC codes: \texttt{[A-Z]\{4\}0[A-Z0-9]\{6\}}, account numbers: 9--18 digits), regex-based approaches combined with contextual heuristics offer higher precision.

Pseudonymization, as defined by the GDPR~\cite{gdpr2016}, replaces identifying data with reversible tokens, allowing data utility while protecting privacy. Our system extends this with AES-256-CBC encryption of the mapping table itself, providing defense-in-depth.

\section{Multi-Tenant Data Isolation}

In multi-tenant systems, data isolation ensures that one tenant's data is inaccessible to others. Johnson et al.~\cite{johnson2023vectordb} discussed metadata-based filtering in vector databases as a lightweight isolation mechanism suitable for cloud-native applications. Our system implements this using Pinecone's native metadata filter with \texttt{owner\_id} tags.

\section{Summary of Literature Gaps}

Table~\ref{tab:lit_gaps} summarizes the security features addressed by existing work versus our proposed system.

\begin{table}[H]
\centering
\caption{Comparison of security features in RAG systems}
\label{tab:lit_gaps}
\begin{tabular}{lcccc}
\toprule
\textbf{Security Feature} & \textbf{Standard RAG} & \textbf{LangChain} & \textbf{LlamaIndex} & \textbf{Ours} \\
\midrule
Prompt Injection Detection & \ding{55} & Partial & Partial & \ding{51} \\
Content Sanitization & \ding{55} & \ding{55} & \ding{55} & \ding{51} \\
PII Detection \& Pseudonymization & \ding{55} & \ding{55} & \ding{55} & \ding{51} \\
AES-256 Encrypted Mappings & \ding{55} & \ding{55} & \ding{55} & \ding{51} \\
Differential Privacy on Embeddings & \ding{55} & \ding{55} & \ding{55} & \ding{51} \\
Multi-Tenant Isolation & \ding{55} & \ding{55} & Partial & \ding{51} \\
RBAC (3-tier) & \ding{55} & \ding{55} & \ding{55} & \ding{51} \\
Controlled De-pseudonymization & \ding{55} & \ding{55} & \ding{55} & \ding{51} \\
\bottomrule
\end{tabular}
\end{table}
